% Source: https://artofproblemsolving.com/community/c3h2050718_fun_problem_i_made
% Stable identity: aops-topic-2050718-q01
Seat the players $A,B,C$ around the circle and let the direction be $d=+1$ (the order $A\to B\to C\to A$) or $d=-1$. With only three players, the player after the next one is the same as the previous one. So whether $X$ plays a skip or a reverse, \emph{the next turn goes to the player before $X$ in the current direction}. The only difference is that a skip keeps the direction and a reverse flips it. Writing states as (player to move, direction):
\begin{itemize}
\item $(A,+)\to(C,+)$ and $(A,-)\to(B,-)$ (Alice skips);
\item $(B,+)\to(A,-)$ and $(B,-)\to(C,+)$ (Bob reverses);
\item $(C,+)\to(B,+)$ or $(B,-)$, and $(C,-)\to(A,+)$ or $(A,-)$, each with probability $\dfrac{1}{2}$ (Carol).
\end{itemize}
No state leads to $(C,-)$, and $(A,+)$ can only be reached from $(C,-)$. So after at most two turns the play stays forever among $(C,+),(B,+),(A,-),(B,-)$, moving by the rules
\[(C,+)\to(B,+)\to(A,-)\to(B,-)\to(C,+)\quad\text{or}\quad(C,+)\to(B,-)\to(C,+).\]
Thus play consists of blocks that start at a Carol turn $(C,+)$: with probability $\dfrac{1}{2}$ the block is $C,B,A,B$ (four turns, one of them Alice's), and with probability $\dfrac{1}{2}$ it is $C,B$ (two turns, none of them Alice's). The blocks are independent, so by the law of large numbers the long-run fraction of Alice's turns is
\[\dfrac{\text{expected Alice turns per block}}{\text{expected turns per block}}=\dfrac{\dfrac{1}{2}\cdot1}{\dfrac{1}{2}\cdot4+\dfrac{1}{2}\cdot2}=\boxed{\dfrac{1}{6}}.\]
(Carol gets $\dfrac{1}{3}$ and Bob $\dfrac{1}{2}$ of the turns.)
